% Informe de la Entrega 2 · Lumina Bank
% Compilar desde doc/latex/:  tectonic informe_entrega2.tex
\documentclass[11pt,a4paper]{article}

\usepackage[spanish,es-nodecimaldot,es-tabla]{babel}
\usepackage{fontspec}
\setmainfont{texgyrepagella}[Extension=.otf,UprightFont=*-regular,BoldFont=*-bold,ItalicFont=*-italic,BoldItalicFont=*-bolditalic]
\setsansfont{texgyreheros}[Extension=.otf,UprightFont=*-regular,BoldFont=*-bold,ItalicFont=*-italic,BoldItalicFont=*-bolditalic]
\setmonofont{texgyrecursor}[Extension=.otf,UprightFont=*-regular,BoldFont=*-bold,Scale=0.9]
\usepackage[margin=2.3cm,top=2.4cm,bottom=2.4cm]{geometry}
\usepackage{graphicx}
\usepackage{booktabs}
\usepackage{xltabular}
\usepackage{array}
\usepackage{float}
\usepackage[section]{placeins}
\usepackage{caption}
\usepackage[table]{xcolor}
\usepackage{etoolbox}
\usepackage{enumitem}
\usepackage{fancyhdr}
\usepackage{titlesec}
\usepackage{needspace}
\usepackage{rotating}
\usepackage[hidelinks]{hyperref}

\definecolor{marino}{HTML}{14213D}
\definecolor{verdeazul}{HTML}{2A9D8F}
\definecolor{gris}{HTML}{6B7385}
\definecolor{fila}{HTML}{F2F5F9}

\captionsetup{font=small,labelfont={bf,color=marino},labelsep=period,justification=centering}
\setlist{itemsep=2pt,topsep=4pt}
\setlength{\parskip}{5pt}
\setlength{\parindent}{0pt}
\setlength{\emergencystretch}{3em}
\widowpenalty=10000
\clubpenalty=10000
\displaywidowpenalty=10000
\renewcommand{\arraystretch}{1.25}
\titleformat{\section}{\Large\sffamily\bfseries\color{marino}}{\thesection}{0.7em}{}
\titleformat{\subsection}{\large\sffamily\bfseries\color{marino}}{\thesubsection}{0.6em}{}
\titleformat{\subsubsection}{\normalsize\sffamily\bfseries\color{verdeazul}}{\thesubsubsection}{0.5em}{}
\titlespacing*{\section}{0pt}{16pt}{6pt}
\titlespacing*{\subsection}{0pt}{12pt}{4pt}

\pagestyle{fancy}
\fancyhf{}
\fancyhead[L]{\small\color{gris}Lumina Bank · Entrega 2}
\fancyhead[R]{\small\color{gris}Taller de Arquitectura Big Data · UTEC}
\fancyfoot[C]{\small\thepage}
\renewcommand{\headrulewidth}{0.3pt}

% Columna X alineada a la izquierda y tablas largas que repiten el encabezado al saltar de página
\newcolumntype{L}{>{\raggedright\arraybackslash}X}
\newcolumntype{P}[1]{>{\raggedright\arraybackslash}p{#1}}
\DeclareRobustCommand{\cod}[1]{{\small\ttfamily\codbreaks #1}}
\DeclareRobustCommand{\codbreaks}{\let\oldus\_\renewcommand{\_}{\oldus\allowbreak}}

% Figura que nunca supera el 80 % de la altura de la página
\newcommand{\figura}[4][0.98]{%
  \begin{figure}[H]\centering
  \includegraphics[width=#1\linewidth,height=0.80\textheight,keepaspectratio]{#2}
  \caption{#3}\label{#4}
  \end{figure}}
% Diagrama ancho: ocupa una página horizontal completa
\newcommand{\figurahorizontal}[3]{%
  \begin{sidewaysfigure}\centering
  \includegraphics[width=\textheight,height=0.9\textwidth,keepaspectratio]{#1}
  \caption{#2}\label{#3}
  \end{sidewaysfigure}}
% Captura de terminal: más baja, para que entren dos por página cuando se puede
\newcommand{\evid}[3]{%
  \begin{figure}[H]\centering
  \includegraphics[width=\linewidth,height=0.42\textheight,keepaspectratio]{evidencias/#1}
  \caption{#2}\label{fig:#3}
  \end{figure}}

\graphicspath{{../img/}}

\begin{document}

% ====================================================================== Portada
\begin{titlepage}
\centering
\vspace*{2.2cm}
{\sffamily\color{gris}\large Universidad de Ingeniería y Tecnología · Maestría\par}
{\sffamily\color{gris}Taller de Arquitectura Big Data · Ciclo 2026-2\par}
\vspace{2.4cm}
{\sffamily\bfseries\color{marino}\fontsize{30}{36}\selectfont Lumina Bank\par}
\vspace{0.4cm}
{\sffamily\color{marino}\Large Plataforma de datos batch y streaming en Google Cloud\par}
\vspace{0.3cm}
{\sffamily\color{verdeazul}\large Entrega 2: procesamiento, transformación y modelado\par}
\vspace{2.6cm}
{\large Bryan Montoya\par David Choque\par Henry Godiño\par}
\vfill
{\color{gris}Lima, 11 de noviembre de 2026\par}
\end{titlepage}

\tableofcontents
\newpage

% ====================================================================== 1
\section{Qué cambiamos después de la Entrega 1}

La revisión de la primera entrega nos dejó cuatro observaciones. Las atendimos todas y, además, desplegamos en un proyecto de Google Cloud (\cod{lumina-bank-utec}) los caminos batch y streaming, la capa Gold, los modelos de BigQuery ML, la API y buena parte de las capas transversales, para medir en lugar de suponer. No es la plataforma completa: la tabla~\ref{tab:desplegado} indica qué quedó fuera y por qué.

El laboratorio corre en \cod{us-central1} y en la multirregión US por costo y porque ahí están disponibles todos los servicios que usamos. El diseño sigue siendo el de la Entrega 1: región principal en Santiago (\cod{southamerica-west1}) y continuidad en São Paulo. Por eso los costos se calculan con precios de \cod{us-central1} y luego se ajustan a Santiago (sección~\ref{sec:costos}). La tabla~\ref{tab:observaciones} resume cada cambio.

\begin{table}[H]
\centering\small
\caption{Observaciones de la docente y respuesta del equipo}
\label{tab:observaciones}
\begin{tabularx}{\linewidth}{P{3.6cm}LP{1.6cm}}
\toprule
\textbf{Observación} & \textbf{Qué hicimos} & \textbf{Sección}\\
\midrule
Mostrar las conexiones entre componentes & Rehicimos la arquitectura en notación C4 con LikeC4: una vista de contexto, una por cada camino de datos y una por cada capa transversal. Cada flecha dice qué dato viaja y por qué canal. Numeramos 18 conexiones y las describimos en una tabla. & \ref{sec:arquitectura}\\
\rowcolor{fila} Incluir gobierno de datos, seguridad y redes & Las tres capas tienen su propia vista C4. En el laboratorio desplegamos la VPC privada, la llave KMS con rotación, la taxonomía de policy tags, el escaneo de calidad, el linaje, las vistas por entidad, la auditoría y Dataflow sin IP pública. Falta asignar los policy tags a las columnas y llevar Dataproc a la VPC privada. VPC Service Controls, Cloud VPN y el balanceador con Cloud Armor quedan diseñados; la tabla~\ref{tab:desplegado} dice por qué. & \ref{sec:transversales}\\
Cloud Run sobredimensionado & Lo recalculamos con la ley de Little sobre el volumen real. El servicio de autorización pasa de 2–100 a 2–20 instancias y la API de datos, de 2–100 a 0–3. & \ref{sec:cloudrun}\\
\rowcolor{fila} Dimensionar y validar Dataflow & Hicimos tres pruebas de carga con un worker fijo. La capacidad real resultó muy por debajo de lo que suponíamos; con ese dato fijamos el rango de workers de la fase MVP y una meta de optimización para producción. & \ref{sec:dataflow}\\
Justificar cada componente & Para cada servicio indicamos qué hace en la solución, por qué lo elegimos y qué alternativa descartamos. & \ref{sec:justificacion}\\
\bottomrule
\end{tabularx}
\end{table}

\FloatBarrier

% ====================================================================== 2
\section{Arquitectura en notación C4}\label{sec:arquitectura}

Usamos el modelo C4 porque separa los niveles de detalle: primero quién usa la plataforma y con qué sistemas se conecta (contexto) y después qué piezas la componen (contenedores). El modelo está escrito en LikeC4 (\cod{doc/arquitectura/lumina.c4}), así que el diagrama se versiona junto con el código y se regenera con un comando. Cada componente lleva el ícono oficial del servicio de Google Cloud que representa.

\subsection{Nivel 1: contexto}

La plataforma recibe transacciones de los canales de las 12 entidades, cambios de saldos de sus cores bancarios, telemetría de cajeros y datos públicos. Entrega alertas al equipo antifraude, tableros a la gerencia y a cada entidad sus propios datos junto con la señal de fraude compartida.

\figura[0.78]{c4/contexto.png}{Diagrama de contexto (C4, nivel 1)}{fig:contexto}

\subsection{Nivel 2: camino streaming}

Es el camino del fraude en tiempo real. La transacción entra por Pub/Sub, Dataflow la evalúa con estado por cliente y por destino, deja las \textit{features} en Memorystore para el servicio de autorización y escribe el resultado en BigQuery. Las alertas salen por otro tópico hacia el SOC. Una copia cruda de cada evento queda en el Data Lake para auditoría y reproceso. El diagrama está en la figura~\ref{fig:streaming} (página~\pageref{fig:streaming}).


\subsection{Nivel 2: camino batch, modelado y exposición}

El histórico y los cierres llegan al Data Lake por Datastream y Storage Transfer. Spark en Dataproc Serverless los limpia y los deja en Parquet; de ahí pasan a Silver y, con SQL, a Gold. BigQuery ML entrena sobre Gold y devuelve el pronóstico de liquidez. Managed Service for Apache Airflow (antes Cloud Composer) orquesta el proceso diario (figura~\ref{fig:batch}, página~\pageref{fig:batch}).


\subsection{Conexiones entre componentes}

La numeración coincide con la de los diagramas.

\begingroup\small
\begin{xltabular}{\linewidth}{@{}P{0.6cm}P{3.9cm}LP{3.6cm}@{}}
\caption{Conexiones de la arquitectura}\label{tab:conexiones}\\
\toprule
\textbf{N.°} & \textbf{Origen → destino} & \textbf{Qué viaja} & \textbf{Canal}\\
\midrule\endfirsthead
\toprule
\textbf{N.°} & \textbf{Origen → destino} & \textbf{Qué viaja} & \textbf{Canal}\\
\midrule\endhead
\bottomrule\endfoot
1 & Canales → Pub/Sub \cod{tx-eventos} & Evento de transacción en JSON & gRPC con TLS; esquema en el tópico solo diseñado\\
2 & Pub/Sub → Dataflow & El mismo evento & Suscripción \textit{pull}, deduplicación por \cod{event\_id}\\
3 & Dataflow → Memorystore & Último score, ubicación y conteo de 24 h por cliente & Redis dentro de la VPC\\
4 & Dataflow → BigQuery Silver & Transacción con score, decisión y latencia & Storage Write API\\
5 & Dataflow → Pub/Sub \cod{alertas-fraude} & Transacciones REVISAR o BLOQUEAR & Publicación en Pub/Sub\\
6 & Pub/Sub → Data Lake & Evento crudo & Suscripción a Cloud Storage cada 5 min\\
7 & Terminales → Pub/Sub → Bronze & Telemetría & Suscripción BigQuery, sin código\\
8 & Cores → Datastream → Data Lake & Cambios de saldos y cierres (CDC) & Cloud VPN de alta disponibilidad\\
9 & Data Lake ↔ Dataproc & CSV crudo de entrada, Parquet curado de salida & Conector de GCS, workers sin IP pública. En el laboratorio el batch corrió en la VPC \cod{default}\\
10 & Data Lake → Silver & \cod{transacciones\_batch} & \cod{bq load} de Parquet\\
11 & Datos públicos → BigQuery SQL & SEC, Google Trends y EDI & Consulta directa y BigQuery sharing\\
12 & Silver y Bronze → SQL → Gold & Dimensiones y hechos & SQL dentro de BigQuery\\
13 & Gold → BigQuery ML & Series de efectivo y transacciones etiquetadas & \cod{CREATE MODEL}, \cod{ML.FORECAST}\\
14 & Memorystore → autorización & \textit{Features} del cliente & Direct VPC egress\\
15 & Silver y Gold → API de datos & Alertas, riesgo, liquidez y KPI & Consultas parametrizadas\\
16 & Gold → Looker Studio & Vistas \cod{v\_kpi\_*} & Conector de BigQuery con BI Engine\\
17 & Gold → 12 entidades & Datos propios y señal de fraude & Vistas autorizadas y filtro por entidad\\
18 & Pub/Sub → SOC & Alertas & Suscripción \textit{push}\\
\end{xltabular}
\endgroup

\figurahorizontal{c4/streaming.png}{Contenedores del camino streaming (C4, nivel 2)}{fig:streaming}
\figurahorizontal{c4/batch.png}{Contenedores del camino batch y de exposición (C4, nivel 2)}{fig:batch}

\subsection{Capas de gobierno, seguridad y redes}\label{sec:transversales}

Cada capa tiene su vista. Arriba están los controles y abajo, los componentes sobre los que actúan.

\figura{c4/gobierno.png}{Capa de gobierno de datos}{fig:gobierno}

\begin{itemize}
\item \textbf{Knowledge Catalog} (antes Dataplex; el comando sigue siendo \cod{gcloud dataplex}) cataloga las tablas, registra el linaje y corre el escaneo de calidad \cod{lumina-dq-\allowbreak transacciones}: unicidad de \cod{event\_id}, nulos en \cod{cliente\_id}, montos no negativos, tipos válidos y distritos entre 1 y 8. Por presupuesto el escaneo evalúa una muestra del 10 \% de \cod{transacciones\_batch}, así que no garantiza la unicidad global de \cod{event\_id}; esa la comprobamos con \cod{COUNT(DISTINCT event\_id)} sobre la tabla completa (figura~\ref{fig:funcional}).
\item \textbf{Policy tags}: se crean en Knowledge Catalog (antes Data Catalog, que sigue siendo el nombre del servicio en IAM y en la API) pero se aplican sobre las columnas de BigQuery. La taxonomía \cod{lumina\_sensibilidad} distingue \cod{PII\_ALTA} (identificadores y saldos) de \cod{PII\_MEDIA} (ubicación y dispositivo). En el laboratorio creamos la taxonomía y los tags (figura~\ref{fig:policytags}), pero todavía no los asignamos a las columnas de BigQuery, así que el control de acceso por columna aún no opera.
\item \textbf{Contratos de datos}: en producción un esquema de Pub/Sub en \cod{tx-eventos} rechazará en la ingesta los mensajes con estructura o tipos inválidos. En el laboratorio ese esquema no está creado: toda la validación la hace la etapa \cod{Parsear} de Dataflow (campos, tipos y montos positivos), que manda lo inválido al tópico \cod{tx-eventos-dlq}. Con el esquema activo, a esa etapa solo le quedarían las reglas de dominio.
\item \textbf{Señal de fraude compartida}: la vista \cod{v\_senal\_fraude\_compartida} publica destinos sospechosos sin datos personales. Responde a la queja del caso de que las entidades no comparten información de fraude.
\end{itemize}

\figura{c4/seguridad.png}{Capa de seguridad}{fig:seguridad}

\begin{itemize}
\item \textbf{IAM}: el diseño pide una cuenta de servicio por componente, con roles sobre el dataset o el bucket que usa. En el laboratorio todos los pipelines comparten \cod{sa-lumina-pipelines}, que tiene roles amplios a nivel de proyecto (\cod{bigquery.dataEditor} y \cod{storage.objectAdmin}, figura~\ref{fig:iam}). Separarla es un pendiente.
\item \textbf{Cloud KMS}: la llave \cod{lumina-datos} rota cada 90 días y hoy cifra el dataset Gold y el bucket de archivo. Silver, Bronze y la zona raw siguen con llaves de Google; extender la llave a esas capas está en los pendientes.
\item \textbf{Datos personales}: la vista \cod{v\_transacciones\_entidad} filtra por la entidad del usuario, reemplaza el id del cliente por un seudónimo y redondea la ubicación.
\item \textbf{Auditoría}: un sink manda los Cloud Audit Logs al dataset \cod{lumina\_auditoria}.
\end{itemize}

\figura{c4/redes.png}{Capa de redes. El perímetro de VPC Service Controls se define en la organización y agrupa proyectos}{fig:redes}

\begin{itemize}
\item \textbf{VPC \cod{lumina-vpc}} con subred privada y Private Google Access. Los workers de Dataflow corren sin IP pública dentro de ella desde la corrida de carga 2 (también la 3). \textit{Desplegado.} Dataproc Serverless todavía no: el job batch se lanzó sin \cod{--subnet} y corrió en la VPC \cod{default}, con IP internas pero fuera de \cod{lumina-vpc}.
\item \textbf{Cloud VPN} lleva el tráfico de Datastream hacia los cores sin pasar por Internet. \textit{Diseñado.}
\item \textbf{VPC Service Controls} no se configura por servicio ni por proyecto sino en la organización: el perímetro agrupa proyectos y bloquea que los datos de BigQuery, Cloud Storage y Pub/Sub salgan de ellos, aunque la credencial sea válida. Planeamos dos perímetros, uno de alta confianza para el proyecto con PII y Gold y otro para el de procesamiento, con reglas de entrada y salida explícitas entre ellos (figura~\ref{fig:redes}). \textit{Diseñado.}
\item \textbf{Cloud Load Balancing y Cloud Armor} son la única entrada pública y hacen el \textit{failover} hacia São Paulo. \textit{Diseñado.} Mientras no existan, la API de datos se publica en su URL de Cloud Run y solo la protege IAM: el tráfico no pasa por el WAF. En producción el servicio se configura con \cod{--ingress internal-and-cloud-load-balancing}, de modo que la única entrada sea el balanceador.
\end{itemize}

\begin{table}[H]
\centering\small
\caption{Qué quedó desplegado en el laboratorio y qué no}
\label{tab:desplegado}
\begin{tabularx}{\linewidth}{P{3.6cm}P{2cm}L}
\toprule
\textbf{Componente} & \textbf{Estado} & \textbf{Motivo}\\
\midrule
Pub/Sub, Dataflow, Dataproc, BigQuery, BigQuery ML, Cloud Run, Cloud Storage & Desplegado & Núcleo de la entrega\\
\rowcolor{fila} VPC, firewall, KMS, auditoría, taxonomía de policy tags, Knowledge Catalog (calidad y linaje) & Desplegado & Capas transversales pedidas en la revisión\\
Asignación de policy tags a columnas & No & Se hace a mano en el esquema de cada tabla; queda para la Entrega 3\\
\rowcolor{fila} Dataproc en \cod{lumina-vpc} y CMEK en el bucket curated & No & El batch corrió en la VPC \cod{default} y curated usa llaves de Google; se corrige al pasar el ETL al DAG\\
Esquema de Pub/Sub en \cod{tx-eventos} & No & La validación la hace Dataflow, con salida a la DLQ\\
Managed Service for Apache Airflow (antes Cloud Composer) & No & Costo: un entorno pequeño cuesta varias veces el presupuesto de S/~180 del laboratorio. Los DAG están escritos y se mostrarán en la Entrega 3\\
\rowcolor{fila} Memorystore & No & Lo consume el servicio de autorización, que todavía no está desplegado; tiene costo fijo por hora\\
VPC Service Controls & No & Los perímetros se definen en una organización de Google Cloud; el proyecto está en una cuenta personal sin organización\\
\rowcolor{fila} Cloud VPN y Datastream & No & No hay cores bancarios reales a los que conectarse; el histórico entra como archivo\\
Load Balancing y Cloud Armor & No & Costo fijo mensual y una sola API de prueba; la API queda protegida solo por IAM\\
\bottomrule
\end{tabularx}
\end{table}

% ====================================================================== 3
\section{Por qué elegimos cada componente}\label{sec:justificacion}

Los criterios fueron cinco: la latencia que pide cada caso de uso, el volumen y su variación, el requisito de código abierto del caso, el trabajo de operación que le exige al equipo y el costo.

\begingroup\small
\begin{xltabular}{\linewidth}{@{}P{2.6cm}LLP{3.5cm}@{}}
\caption{Justificación técnica por componente}\label{tab:justificacion}\\
\toprule
\textbf{Componente} & \textbf{Qué hace en Lumina} & \textbf{Por qué este} & \textbf{Descartamos}\\
\midrule\endfirsthead
\toprule
\textbf{Componente} & \textbf{Qué hace en Lumina} & \textbf{Por qué este} & \textbf{Descartamos}\\
\midrule\endhead
\bottomrule\endfoot
\multicolumn{4}{@{}l}{\textbf{\color{verdeazul}Ingesta}}\\
Pub/Sub & Recibe transacciones y telemetría y las retiene 7 días & Escala de 175 a 5\,000 TPS sin dimensionar particiones. Tiene suscripciones nativas a BigQuery y Cloud Storage que nos ahorran dos pipelines & Kafka gestionado: hay que dimensionar brokers y particiones\\
Datastream & Copia los cambios de los cores & CDC sin instalar agentes en las bases de las entidades & Debezium: hay que operar Kafka Connect\\
Storage Transfer & Carga PaySim, feriados y archivos de las entidades & Programado, con reintentos y verificación & Scripts en una VM\\
BigQuery sharing (antes Analytics Hub) & Lee EDI y publica la señal de fraude & Consume datos del proveedor sin copiarlos y comparte datos entre entidades con permisos & Copiar archivos\\
\multicolumn{4}{@{}l}{\textbf{\color{verdeazul}Procesamiento}}\\
Dataflow (Apache Beam) & Evalúa el fraude en tiempo real & Estado y \textit{timers} gestionados para viaje imposible y cuenta mula; escala horizontal. Beam también corre en Flink o Spark & Spark Streaming: micro-lotes y clúster permanente\\
Dataproc Serverless (Spark) & ETL batch de CSV a Parquet & Spark es el estándar y lo conocemos; se paga por job & Clúster fijo: se paga aunque no corra\\
BigQuery SQL & Construye Gold y cruza con datos públicos & Transforma donde está el dato, sin moverlo & Hacer todo en Spark\\
BigQuery ML & Fraude (XGBoost) y liquidez (ARIMA\_PLUS) & Entrena con SQL sobre los mismos datos & Gemini Agent Platform (antes Vertex AI) desde el inicio: más trabajo para el MVP\\
Managed Service for Apache Airflow & Orquesta el batch y el precalentamiento & Airflow: dependencias, reintentos y \textit{backfills} & Workflows: no hace \textit{backfill}\\
\multicolumn{4}{@{}l}{\textbf{\color{verdeazul}Almacenamiento}}\\
Memorystore & \textit{Features} para autorizar en menos de 150 ms & Lecturas de submilisegundos; Valkey y Redis son abiertos & Bigtable: más latencia y mayor costo mínimo\\
Cloud Storage & Data Lake con zonas raw, curated y archive & Barato y durable, con ciclo de vida y bloqueo WORM; formatos abiertos & HDFS: une cómputo y almacenamiento\\
BigQuery & Bronze, Silver, Gold y modelos & Serverless, con seguridad por fila y columna y acceso directo a SEC y Trends & Hive o PostgreSQL: operación o escala\\
\multicolumn{4}{@{}l}{\textbf{\color{verdeazul}Exposición}}\\
Cloud Run & Autorización y API de datos & Contenedores que escalan por concurrencia; la API escala a cero & GKE: un clúster que operar\\
Looker Studio & Tableros de gerencia & Sin licencia; con BI Engine responde en menos de un segundo & Looker: licencia\\
Vistas autorizadas & Datos por entidad & Control nativo sin copiar datos por banco & Un dataset por entidad\\
\multicolumn{4}{@{}l}{\textbf{\color{verdeazul}Capas transversales}}\\
Knowledge Catalog & Catálogo, linaje y calidad & Integrado con BigQuery y Dataflow sin instalar nada & DataHub o Atlas: hay que desplegarlos\\
Cloud KMS & Llaves propias con rotación & Permite revocar el acceso y demostrar que la entidad controla las llaves de sus datos en la nube & Llaves gestionadas por Google\\
VPC Service Controls & Perímetro de datos a nivel de organización & Agrupa proyectos e impide copiar datos fuera aunque la credencial sea válida & Solo IAM\\
Load Balancing + Armor & Entrada pública y \textit{failover} & Balanceo global entre Santiago y São Paulo, con WAF & Balanceo regional\\
\end{xltabular}
\endgroup

% ====================================================================== 4
\section{Dimensionamiento de Dataflow y Cloud Run}

\subsection{Volumen de referencia}

\begin{table}[H]
\centering\small
\caption{Escenarios de carga}
\label{tab:escenarios}
\begin{tabularx}{0.85\linewidth}{LP{5.5cm}r}
\toprule
\textbf{Escenario} & \textbf{Cálculo} & \textbf{TPS}\\
\midrule
Promedio diario & 15 M de transacciones ÷ 86\,400 s & 175\\
\rowcolor{fila} Hora pico de un día normal & 2,5 × el promedio & 440\\
Día de pago o cierre de mes & 60 M × 12 \% en la hora pico ÷ 3\,600 s & 2\,000\\
\rowcolor{fila} Pico de diseño (Cyber Days) & 2,5 × el día de pago & 5\,000\\
\bottomrule
\end{tabularx}
\end{table}

En la Entrega 1 sumamos la telemetría y el \textit{clickstream} a la carga de Dataflow. Fue un error: la telemetría entra a BigQuery por una suscripción directa y el \textit{clickstream} no pasa por el pipeline de fraude. Por Dataflow solo pasan transacciones.

\subsection{Dataflow: lo que medimos}\label{sec:dataflow}

Al empezar esta entrega estimamos, sin medir, que un worker \cod{n2-standard-2} procesaría unas 1\,000 transacciones por segundo. Con eso bastaban 8 workers para el pico de diseño. Antes de fijar ese número lo pusimos a prueba.

\textbf{Cómo medimos.} Dejamos el job fijo en un solo worker (\cod{--max\_num\_workers 1}) y publicamos transacciones de PaySim por escalones de tres minutos: 100, 300, 600 y 900 por segundo. En cada minuto registramos cuántos eventos llegaron a BigQuery, la latencia p95 entre el envío y la decisión, el backlog de la suscripción y el uso de CPU del worker. La consulta está en \cod{src/sql/06\_validacion\_capacidad.sql}. Solo la corrida 1 recorrió el protocolo completo de 12 minutos. Como el worker se saturó ya en el primer escalón, las corridas 2 y 3 buscaban solo ver si el tope de 100 eventos por segundo cambiaba con cada ajuste, y las detuvimos antes: la 2 a los 6 minutos y la 3 a los 4. Por eso en la figura~\ref{fig:carga} se comparan los eventos por segundo sostenidos; la latencia p95 crece con el backlog acumulado y no es comparable entre corridas de distinta duración.

\textbf{Qué encontramos.} El worker se saturó en el primer escalón: le llegaban 100 eventos por segundo y procesaba entre 65 y 75. En la corrida 1 el simulador publicó 342\,000 eventos en los cuatro escalones; el worker procesó 62\,159 antes de que canceláramos el job y el backlog seguía creciendo (162\,392 mensajes a las 04:37, figura~\ref{fig:backlog}). Hicimos tres corridas de carga (tabla~\ref{tab:corridas} y figura~\ref{fig:carga}). La 1 y la 3 cambian una sola cosa; la 2 cambió dos a la vez, los hilos y la red, así que su resultado no se puede atribuir a ninguna de las dos por separado.

\begin{table}[H]
\centering\small
\caption{Corridas de la prueba de carga con un worker \cod{n2-standard-2}}
\label{tab:corridas}
\begin{tabularx}{\linewidth}{P{1.3cm}P{3.5cm}P{1.9cm}P{1.6cm}L}
\toprule
\textbf{Corrida} & \textbf{Cambio} & \textbf{ev/s sostenidos} & \textbf{CPU} & \textbf{Qué aprendimos}\\
\midrule
1 & Configuración por defecto & 65–75 (pico de 143) & 93–97 \% & El worker no sigue el ritmo: el backlog llegó a 162\,000 mensajes y la latencia p95, a casi 9 minutos. El runner advirtió que la caché de estado era insuficiente y sugirió \cod{--max\_cache\_memory\_usage\_mb}.\\
\rowcolor{fila} 2 & 64 hilos por proceso y, a la vez, workers sin IP pública en la VPC privada & 55–63 & 94–96 \% & No mejoró. La CPU al 95 \% apunta a un límite de cómputo, pero no probamos la sugerencia del runner sobre la caché de estado, así que no podemos descartar que también influya. No medimos por separado el efecto de la red; es la configuración que usará producción, así que la tomamos como piso.\\
3 & Escritura a BigQuery en lotes (\cod{with\_auto\_sharding}), con la misma red de la 2 (sin IP pública en \cod{lumina-vpc}) & 78–89 & — & Mejora cerca de un 15 \%. El sink no era el problema principal.\\
\bottomrule
\end{tabularx}
\end{table}

\figura{carga.png}{Eventos procesados y latencia p95 por minuto en las tres corridas}{fig:carga}

Con poca carga el pipeline sí cumple: a 5 eventos por segundo la latencia p95 fue de 2,4 s, por debajo de la meta de 5 s. El problema aparece al saturarse. Haciendo la cuenta (2 vCPU al 95 \% para unos 70 eventos por segundo), cada transacción consume unos 27 ms de CPU. El costo está en las dos etapas con estado: por cada evento hay una lectura, una escritura y un \textit{timer} por cliente y lo mismo por destino, y en el SDK de Python cada una de esas operaciones se serializa y viaja entre el proceso de Python y el \textit{runner}. Escribimos ese código pensando en que fuera claro, no en que fuera rápido, y la prueba lo dejó a la vista.

\textbf{Qué significa para el dimensionamiento.} Con 75 eventos por segundo por worker, usados al 70 \%, el código actual necesita 4 workers en promedio, 39 en un día de pago y 96 en el pico de diseño. Pero dentro de la VPC privada, que es como correrá en producción, el worker sostuvo entre 55 y 63. Para cubrir el día de pago con 40 workers haría falta que cada uno procese al menos 71,4 eventos por segundo ($2\,000 / (40 \times 0{,}7)$), y la corrida 2 no llega. Con el piso de 55, el día de pago pide 52 workers. Con la meta de optimización que fijamos más abajo, harían falta 1, 6 y 15 (figura~\ref{fig:dimensionamiento}).

Estas cifras suponen que la capacidad crece en línea recta con el número de workers. No lo demostramos: medimos con un solo worker. En Dataflow el estado se reparte por clave entre los workers y con 6,9 millones de clientes la carga queda bien distribuida, así que esperamos un comportamiento cercano al lineal, pero el \textit{shuffle} y las escrituras a BigQuery pueden restar eficiencia. La prueba con autoscaling de la Entrega 3 lo confirmará.

\figura{dimensionamiento.png}{Workers de Dataflow e instancias de Cloud Run por escenario}{fig:dimensionamiento}

Esto cambia la lectura de la Entrega 1. Allí el máximo de 30 workers \cod{n2-standard-4} no salió de una medición: sobrestimaba la carga (contaba la telemetría) y suponía una capacidad por worker que el código no alcanza. Ahora los dos datos están medidos o corregidos, y la decisión es esta:

\begin{enumerate}
\item \textbf{Fase MVP, con el código actual}: mínimo de 4 workers, máximo de 55 y precalentamiento a 20 en días de pago con el DAG \cod{lumina\_precalentamiento\_dataflow} (\cod{src/orchestration/\allowbreak dag\_lumina\_precalentamiento.py}). El máximo sale del peor caso medido (52 workers con 55 eventos por segundo) más un margen pequeño, y cubre hasta el día de pago. Subirlo de 40 a 55 solo cuesta en las horas pico de esos días, unos USD 50 al mes, que caben en la contingencia del 15 \% de la tabla~\ref{tab:costos}. Si llegara un pico mayor, Pub/Sub retiene los mensajes sin perderlos, pero las alertas saldrían con minutos de retraso. Por eso esta versión no debe operar en un Cyber Day.
\item \textbf{Antes de producción}: optimizar el camino con estado hasta al menos 500 eventos por segundo por worker y repetir esta misma prueba. Las opciones, en el orden en que las probaremos: reescribir las dos etapas con estado en Beam para Java (el grafo y el diseño no cambian, y Beam corre igual en Dataflow); quitar el \textit{timer} que se reprograma en cada evento y usar expiración de estado; y leer el historial del destino desde Memorystore en lugar de mantenerlo en el pipeline. Con 500 eventos por segundo, el máximo queda en 15 workers.
\end{enumerate}

En el diseño, la autorización de cada transacción no depende de esto: la decide Cloud Run con las \textit{features} de Memorystore en milisegundos, y Dataflow solo mantiene esas \textit{features} al día y genera las alertas. Un atraso de Dataflow degradaría la calidad de la detección pero no detendría los pagos. Hoy esa protección no existe: ni el servicio de autorización ni Memorystore están desplegados, y el pipeline todavía no escribe en Memorystore (tabla~\ref{tab:beam}). Hasta que lo estén, un atraso de Dataflow atrasa las alertas y no hay una segunda línea que decida en tiempo real.


\subsection{Cloud Run}\label{sec:cloudrun}

Usamos la ley de Little: solicitudes simultáneas = tasa de llegada × tiempo de respuesta. Las instancias quedan limitadas por la CPU o por la concurrencia, lo que llegue primero, y el autoscaler de Cloud Run apunta a un 60 \% de uso.

\textbf{Servicio de autorización.} Hay una solicitud por transacción. Cada una tarda unos 25 ms (una lectura en Memorystore dentro de la VPC más las reglas) y consume unos 2 ms de CPU, así que una instancia de 1 vCPU atiende unas 500 solicitudes por segundo; al 60 \% de uso quedan 300. En el pico de diseño hacen falta $\lceil 5\,000/300 \rceil = 17$ instancias y en un día de pago, 7. Fijamos un máximo de 20 (17 más un margen del 15 \%) y un mínimo de 2, para no tener arranques en frío y seguir funcionando si cae una zona. La concurrencia baja de 80 a 40: con 25 ms por solicitud cada instancia tiene unas 8 en curso, y 80 no aportaba nada.

\textbf{API de datos.} La usan el SOC y las aplicaciones para consultas de riesgo; los tableros van directo a BigQuery. Estimamos 20 solicitudes por segundo en el pico y 1,5 s por consulta, o sea 30 solicitudes simultáneas. Como cada solicitud pasa casi todo el tiempo esperando a BigQuery, una instancia aguanta 20 en paralelo. Con 2 instancias basta; dejamos un máximo de 3, que además limita a 60 las consultas simultáneas contra BigQuery. El mínimo es 0 en el laboratorio y 1 en producción.

\begin{table}[H]
\centering\small
\caption{Configuración de Cloud Run antes y después}
\label{tab:cloudrun}
\begin{tabularx}{\linewidth}{lLLL}
\toprule
\textbf{Parámetro} & \textbf{Entrega 1} & \textbf{Autorización} & \textbf{API de datos}\\
\midrule
Instancias mínimas & 2 & 2 & 0 (laboratorio), 1 (producción)\\
\rowcolor{fila} Instancias máximas & 100 & 20 & 3\\
Concurrencia & 80 & 40 & 20\\
\rowcolor{fila} CPU y memoria & — & 1 vCPU, 512 MiB, CPU siempre asignada & 1 vCPU, 512 MiB, CPU solo durante la solicitud\\
\bottomrule
\end{tabularx}
\end{table}

Bajar el máximo no reduce el costo por sí solo, porque Cloud Run cobra por uso. Lo que hace es poner un techo al gasto y proteger a Memorystore y BigQuery de una avalancha de solicitudes. La API ya está desplegada con esta configuración (figura~\ref{fig:cloudrun}).

\subsection{Impacto en el costo}\label{sec:costos}

Recalculamos el costo mensual con precios de lista de \cod{us-central1}, igual que en la Entrega 1, y lo llevamos a Santiago con un ajuste del 22 \%. La tabla~\ref{tab:costos} muestra solo las partidas que cambian; el resto (Pub/Sub, Dataproc, BigQuery, Memorystore, Cloud Storage, Airflow gestionado, seguridad, observabilidad y DR) suma USD 2\,635 al mes y no se modifica.

\begin{table}[H]
\centering\small
\caption{Costo mensual estimado (USD)}
\label{tab:costos}
\begin{tabularx}{\linewidth}{LP{2.2cm}P{2.6cm}P{2.6cm}}
\toprule
\textbf{Partida} & \textbf{Entrega 1} & \textbf{MVP (código actual)} & \textbf{Producción (optimizado)}\\
\midrule
Dataflow & 1\,100 & 800 & 250\\
\rowcolor{fila} Cloud Run, balanceador y Cloud Armor & 80 & 160 & 160\\
Redes (VPN hacia los cores, NAT, egress) & 80 & 180 & 180\\
\rowcolor{fila} Resto de servicios & 2\,635 & 2\,635 & 2\,635\\
\midrule
Subtotal más 15 \% de contingencia & 4\,480 & 4\,340 & 3\,710\\
\rowcolor{fila} \textbf{Santiago con DR en São Paulo} & \textbf{5\,500} & \textbf{5\,300} & \textbf{4\,500}\\
\textbf{Año 1} & \textbf{66\,000} & \textbf{63\,600} & \textbf{54\,000}\\
\rowcolor{fila} USD por millón de transacciones & 12 & 12 & 10\\
\bottomrule
\end{tabularx}
\end{table}

Dataflow en el MVP supone unos 4,5 workers en promedio a lo largo del día, más los picos de los días de pago. Cloud Run sube respecto de la Entrega 1 porque ahora contamos bien las dos instancias mínimas del servicio de autorización, y la partida de redes crece porque incluye la VPN hacia los cores. Son estimaciones que hay que confirmar en la calculadora de precios de Google Cloud.

El laboratorio corre bajo un presupuesto de S/~180 con alertas al 50, 90 y 100 \% (figura~\ref{fig:presupuesto}). Dataflow se apagó al terminar cada prueba; la API escala a cero cuando nadie la usa.


% ====================================================================== 5
\section{Implementación de la Entrega 2}

\subsection{Pipeline batch con Spark}

El script \cod{src/dataproc/batch\_paysim\_etl.py} lee el CSV de PaySim con un esquema explícito y escribe Parquet particionado por fecha. Descarta filas sin llaves o monto, tipos fuera de dominio, montos negativos y duplicados. Genera un \cod{event\_id} determinista, convierte el paso de simulación en fecha real y asigna distrito, entidad, edad, canal y ubicación con el mismo hash MD5 que usan el simulador y el SQL. También calcula los descuadres de saldo, que son la señal más fuerte de fraude en PaySim.

El job de Dataproc Serverless tardó 2 min 51 s, leyó 6\,362\,620 filas y las escribió todas en 31 particiones (552 MiB en Parquet). Ninguna fila cayó en las reglas de descarte del ETL, pero eso no significa que el archivo viniera limpio: tuvimos que normalizarlo a mano antes de subirlo, fuera del pipeline, como se explica a continuación. Después, \cod{bq load} llevó el Parquet a \cod{lumina\_silver.transacciones\_batch}, particionada por día y agrupada por distrito, tipo y cliente.

Al descargar el dataset de Kaggle encontramos dos problemas que obligaron a normalizarlo antes de subirlo: viene partido en \cod{train.csv} y \cod{validation.csv}, con las columnas \cod{isFraud} e \cod{isFlaggedFraud} en orden invertido, y con fines de línea de Windows. Como el ETL lee por posición, sin esa limpieza la etiqueta de fraude habría quedado cambiada. Esa normalización la hicimos con un script aparte; en la Entrega 3 la pasaremos al ETL, que leerá las columnas por nombre.

\subsection{Pipeline streaming con Apache Beam}

\begin{table}[H]
\centering\small
\caption{Pasos del pipeline \cod{fraud\_streaming\_pipeline.py}}
\label{tab:beam}
\begin{tabularx}{\linewidth}{P{3cm}L}
\toprule
\textbf{Paso} & \textbf{Qué hace}\\
\midrule
\cod{LeerPubSub} & Lee la suscripción y descarta duplicados por \cod{event\_id}\\
\rowcolor{fila} \cod{Parsear} & Valida campos, tipos y montos; lo inválido va al tópico \cod{tx-eventos-dlq}\\
\cod{ViajeImposible} & Guarda por cliente la última ubicación y calcula la velocidad\\
\rowcolor{fila} \cod{VelocidadDestino} & Cuenta por destino las transacciones de los últimos 10 minutos\\
\cod{Reglas} & Suma los pesos de las reglas de la tabla~\ref{tab:reglas} en un score de 0 a 1: desde 0,8 bloquea, desde 0,5 manda a revisión\\
\rowcolor{fila} Salidas & Todas las transacciones a Silver; las alertas a Silver y al tópico \cod{alertas-fraude}. La escritura de \textit{features} a Memorystore está en el diseño, no en el código actual\\
\bottomrule
\end{tabularx}
\end{table}

En la prueba funcional el simulador publicó 30\,000 transacciones a 100 por segundo y las 30\,000 llegaron a Silver con \cod{event\_id} distintos; 130 terminaron en alerta (50 bloqueos y 80 a revisión, figura~\ref{fig:funcional}). La figura~\ref{fig:streamingres} muestra 29\,984 y 122 porque filtra por la hora del evento y deja fuera los últimos segundos de la corrida. En paralelo, 50 terminales enviaron 6\,000 lecturas de telemetría que la suscripción escribió en Bronze sin pasar por ningún pipeline.

\begin{table}[H]
\centering\small
\caption{Reglas de fraude del pipeline y su peso en el score}
\label{tab:reglas}
\begin{tabularx}{\linewidth}{P{4.4cm}LP{1.2cm}}
\toprule
\textbf{Regla} & \textbf{Se activa cuando} & \textbf{Peso}\\
\midrule
\cod{R1\_VACIADO\_CUENTA} & TRANSFER o CASH\_OUT por todo el saldo de origen (la cuenta queda en cero) & 0,5\\
\rowcolor{fila} \cod{R2\_INCONSISTENCIA\_DESTINO} & TRANSFER o CASH\_OUT hacia un destino con saldo cero antes y después & 0,3\\
\cod{R3\_MONTO\_ALTO} & TRANSFER de más de 200\,000 & 0,2\\
\rowcolor{fila} \cod{R4\_VIAJE\_IMPOSIBLE} & Entre dos transacciones del mismo cliente la velocidad supera 800 km/h & 0,5\\
\cod{R5\_CUENTA\_MULA} & Tres o más TRANSFER al mismo destino en 10 minutos & 0,3\\
\bottomrule
\end{tabularx}
\end{table}

\subsection{Batch y streaming en el mismo modelo}

Los dos caminos escriben en Silver con el mismo esquema y asignan distrito y entidad con la misma función. La vista \cod{fact\_transacciones} los une con una columna \cod{origen}, de modo que un KPI calculado sobre el histórico se puede comparar con uno calculado en vivo. Revisamos que la unión no duplique transacciones: ningún \cod{event\_id} aparece en los dos orígenes, y el histórico cubre julio de 2026 mientras que el stream es del 10 de noviembre, así que tampoco se cruzan en los KPI diarios (figura~\ref{fig:funcional}).

\subsection{Modelo estrella en Gold}

\begin{table}[H]
\centering\small
\caption{Tablas y vistas de la capa Gold}
\label{tab:gold}
\begin{tabularx}{\linewidth}{P{4.2cm}P{3.3cm}L}
\toprule
\textbf{Objeto} & \textbf{Grano} & \textbf{Uso}\\
\midrule
\cod{dim\_distrito} & 8 distritos & Liquidez y congestión; cada uno se asocia a una región de Google Trends\\
\rowcolor{fila} \cod{dim\_entidad} & 12 entidades & Acceso por entidad; cada una tiene un banco de la SEC como referencia\\
\cod{dim\_cliente} & 6,9 M de clientes & Segmentación por edad, distrito y entidad\\
\rowcolor{fila} \cod{dim\_fecha} & Día & Feriados y días de pago\\
\cod{fact\_transacciones} & Transacción & Une batch y streaming\\
\rowcolor{fila} \cod{fact\_liquidez\_hora} & Distrito × hora & Entrada del pronóstico de efectivo\\
\cod{fact\_tendencias\_semana} & Distrito × semana & Demanda digital según Google Trends\\
\rowcolor{fila} \cod{fact\_benchmark\_sector} & Banco × año & ROA y apalancamiento del sector\\
\cod{v\_kpi\_*} y vistas de gobierno & — & Informes operativos, reglas de fraude, acceso por entidad\\
\bottomrule
\end{tabularx}
\end{table}

Dos ajustes respecto a lo planeado. Google Trends no publica datos de Perú en \cod{international\_top\_terms}, así que usamos Colombia como referencia andina (33 regiones). Y para elegir los bancos de referencia usamos los informes 10-K de 2018, el último año con cobertura completa en el dataset de la SEC; con 2020 solo aparecían 4 bancos.

\subsection{Modelos en BigQuery ML}

El modelo de fraude (XGBoost) se evaluó sobre el 20 \% de TRANSFER y CASH\_OUT que no vio en el entrenamiento, con la proporción real de fraude. Obtuvo una precisión de 0,81, un recall de 0,995 y un AUC de 0,998, por encima de la meta de la Entrega 1 (recall de al menos 85 \% con precisión de al menos 70 \%). Ese resultado tiene una fuga de información. Las dos variables que más pesan son el descuadre del saldo de origen (\cod{err\_saldo\_orig}, 2,31) y el saldo de origen después de la transacción (\cod{saldo\_nuevo\_orig}, 0,44), y ninguna existe en el momento de autorizar: se conocen cuando la transacción ya se ejecutó. Entrenamos un segundo modelo solo con lo que se sabe antes de decidir (tipo, monto, saldos previos, proporción del saldo que se mueve y hora). Obtuvo una precisión de 0,81, un recall de 0,994 y un AUC de 0,998 (figura~\ref{fig:sinfuga}), prácticamente lo mismo. La variable que manda ahora es la proporción del saldo que se mueve: en PaySim el fraude casi siempre vacía la cuenta, el mismo patrón que usa la regla R1. Con datos reales no esperamos una separación tan limpia; la cifra sirve para validar el flujo, no para prometer ese desempeño.

Las reglas del pipeline, sin modelo, detectan el 98,9 \% de los fraudes con una precisión de 0,52: por cada fraude detectado generan 0,9 falsas alarmas. El modelo genera 0,23, cerca de un 74 \% menos. La comparación es aproximada porque las bases no son iguales. Las reglas se midieron sobre las 92\,074 transacciones reales de PaySim que habían pasado por el pipeline cuando tomamos la captura (figura~\ref{fig:reglas}): las 30\,000 de la prueba funcional y las 62\,159 de la corrida de carga 1, menos 85 eventos sintéticos. La vista excluye los sintéticos, que son los PAYMENT que el simulador inyecta antes de un viaje imposible. La prueba funcional se publicó con \cod{--solo-transfer-cashout}; las corridas de carga, sin ese filtro, así que traen todos los tipos. Sobre las 125\,274 transacciones reales que acumuló el stream al cierre, 73\,120 son TRANSFER o CASH\_OUT (14\,135 y 58\,985) y 52\,154 son PAYMENT, CASH\_IN o DEBIT (31\,212, 20\,335 y 607). Las reglas R1 a R5 solo se activan en TRANSFER y CASH\_OUT, y en PaySim el fraude solo ocurre en esos dos tipos, así que los otros tres suman únicamente verdaderos negativos: inflan \cod{vn}, pero no cambian la precisión ni el recall. El modelo se midió sobre el 20 \% reservado de TRANSFER y CASH\_OUT del histórico.

Los conteos del camino streaming cambian entre figuras porque \cod{transacciones\_stream} creció mientras capturábamos: 92\,074 filas evaluables en la figura~\ref{fig:reglas} (23:38 del 9 de noviembre), 97\,708 en la figura~\ref{fig:gold} (23:44, con la corrida 2 en curso) y 125\,274 en la figura~\ref{fig:funcional}, que tomamos el 10 de noviembre con las tres corridas de carga ya cargadas. Las vistas de Gold agregan sobre la misma tabla; la diferencia es solo el momento del corte.

El pronóstico de liquidez (ARIMA\_PLUS por distrito) no alcanza la meta de un error menor al 15 \%. Con los cuatro días de evaluación, el MAPE diario va de 70 \% a 816 \% según el distrito. Si se excluye el 31 de julio baja a entre 32 \% y 59 \% en siete distritos y queda en 228 \% en el octavo (figura~\ref{fig:mape}). Ninguna de las dos lecturas cumple la meta. La causa está en los datos: en sus últimos días PaySim casi no tiene movimiento (282 transacciones el 31 de julio frente a unas 50\,000 en un día típico), y esos son justamente los días reservados para evaluar. Con un año de datos reales el modelo podrá aprender la estacionalidad semanal y mensual; con 27 días no puede.

La vista \cod{v\_plan\_liquidez\_t1} sumaba las cotas superiores horarias y eso inflaba el efectivo recomendado. La corregimos: ahora devuelve solo T+1 y agrega bien el margen diario (figura~\ref{fig:liquidez}). La precisión del pronóstico queda limitada por los 27 días de PaySim y no es el foco de esta entrega.

\subsection{API de datos}

La API en Cloud Run responde 403 a quien llama sin credenciales y devuelve las alertas, el riesgo de un cliente y el plan de liquidez a quien presenta un token válido (figura~\ref{fig:api}). Hoy se accede directo por su URL de Cloud Run, sin pasar por Cloud Armor (sección~\ref{sec:transversales}). Tiene además dos problemas de gobierno que corregiremos: corre con la cuenta compartida \cod{sa-lumina-pipelines}, que tiene permisos de escritura en todo el proyecto, y consulta Silver directamente, así que devuelve \cod{cliente\_id} en claro (figura~\ref{fig:api}) en lugar de pasar por las vistas seudonimizadas. Debe usar una cuenta propia con lectura solo sobre esas vistas. En la Entrega 3 la conectaremos al tablero y a la demo.

% ====================================================================== 6
\section{Evidencias de la implementación}

Todas las capturas son salidas reales de \cod{gcloud} y \cod{bq} ejecutadas sobre el proyecto \cod{lumina-bank-utec} el 9 y 10 de noviembre de 2026. Cada una muestra el comando, la fecha y el resultado. Los logs completos están en \cod{doc/evidencias/logs/}. Las figuras marcadas como ``Consola'' son capturas de la consola web de Google Cloud y complementan la salida de terminal del mismo recurso.

\subsection{Proyecto e infraestructura base}

\evid{02_presupuesto.png}{Presupuesto de S/~180 con alertas al 50, 90 y 100 \%. La cuenta de facturación está en soles}{presupuesto}
\evid{03_apis.png}{APIs habilitadas en el proyecto}{apis}
\evid{04_setup.png}{Ejecución de \cod{infra/setup.sh}: APIs, buckets, cuenta de servicio y datasets}{setup}
\evid{05_buckets.png}{Los cuatro buckets del Data Lake en la multirregión US; \cod{raw} tiene versionado}{buckets}
\evid{06_ciclo_vida.png}{Ciclo de vida del bucket de archivo: Nearline a 30 días, Coldline a 90 y Archive a 365}{ciclovida}
\evid{07_pubsub.png}{Tópicos y suscripciones de Pub/Sub: Dataflow, archivo en Cloud Storage y escritura directa a BigQuery}{pubsub}

\subsection{Gobierno, seguridad y redes}

\evid{08_gobierno_script.png}{Segunda ejecución de \cod{infra/gobierno\_seguridad\_red.sh}: el script detecta lo que ya creó y no lo duplica}{scriptgob}
\evid{09_red.png}{Subred privada con Private Google Access y reglas de firewall de \cod{lumina-vpc}}{red}
\evid{consola/vpc.png}{Consola: subred \cod{lumina-datos-us-central1} (10.10.0.0/20) con Private Google Access activo y flow logs}{c-vpc}
\evid{10_kms.png}{Llave \cod{lumina-datos}: rotación cada 90 días y permisos solo para los agentes de BigQuery y Cloud Storage}{kms}
\evid{consola/kms.png}{Consola: llave \cod{lumina-datos} disponible en el keyring \cod{lumina-kr}, con la próxima rotación programada para el 7 de febrero de 2027 a las 23:25 en hora de Lima (UTC-5); la CLI de la figura~\ref{fig:kms} muestra la misma fecha en UTC, 8 de febrero a las 04:25}{c-kms}
\evid{11_cmek.png}{El dataset Gold y el bucket de archivo usan la llave propia (CMEK)}{cmek}
\evid{12_iam.png}{Roles de la cuenta de servicio de los pipelines}{iam}
\evid{13_policy_tags.png}{Taxonomía de sensibilidad con los tags \cod{PII\_ALTA} y \cod{PII\_MEDIA}}{policytags}
\evid{14_dataplex_dq.png}{Escaneo de calidad de Knowledge Catalog (\cod{gcloud dataplex}) sobre Silver: cinco reglas aprobadas en una muestra de 636\,421 filas}{dq}
\evid{consola/dataplex_dq.png}{Consola de Knowledge Catalog: dos ejecuciones del escaneo aprobadas y las cinco reglas con umbral de 100 \%}{c-dq}
\evid{16_vista_entidad.png}{La vista por entidad solo devuelve filas de la entidad 3 y reemplaza el id del cliente por un seudónimo}{vistaentidad}
\evid{17_auditoria.png}{Sink de auditoría y tablas de Cloud Audit Logs en BigQuery}{auditoria}
\evid{consola/linaje.png}{Linaje de \cod{fact\_liquidez\_hora} en BigQuery: viene de \cod{fact\_transacciones}, \cod{dim\_distrito} (Google Trends) y \cod{dim\_fecha} (feriados). \cod{dim\_distrito} entra solo como la lista de los ocho distritos sobre la que se arma la serie por hora; Trends le da el nombre de región y no es variable del pronóstico. El tramo de Spark lo publica Dataproc; los pasos de BigQuery los registramos con eventos OpenLineage (\cod{infra/registrar\_linaje.py})}{c-linaje}

\subsection{Pipeline batch}

\evid{18_paysim_raw.png}{PaySim normalizado en la zona raw del Data Lake (469,72 MiB)}{paysimraw}
\evid{19_dataproc_submit.png}{Envío del job de Spark a Dataproc Serverless con la línea de control de calidad}{dataproc}
\evid{19b_dataproc_estado.png}{Estado final del job de Spark: SUCCEEDED en 2 min 51 s}{dataprocestado}
\evid{consola/dataproc_batch.png}{Consola de Dataproc Serverless, que la consola ya muestra como Managed Apache Spark: batch \cod{lumina-etl-20261110-v2} terminado en 2 min 50 s con la cuenta de servicio de los pipelines (0,37 DCU-hora)}{c-dataproc}
\evid{20_curated.png}{Zona curated: 31 particiones por fecha en Parquet}{curated}
\evid{20a_dimensiones.png}{Dimensiones conformadas: distritos asociados a regiones de Trends y entidades con su banco de referencia en la SEC}{dimensiones}
\evid{21_silver_batch.png}{Tabla Silver de transacciones batch: 6\,362\,620 filas, particionada y agrupada}{silver}

\subsection{Pipeline streaming}

\evid{24_simulador.png}{Simulador publicando 30\,000 transacciones a 100 por segundo}{simulador}
\evid{25b_corrida_funcional.png}{Decisiones de las 30\,000 transacciones de la prueba funcional y revisión de duplicados entre el histórico y el stream}{funcional}
\evid{25_streaming.png}{Decisiones del pipeline, alertas con las reglas que las dispararon y telemetría escrita en Bronze}{streamingres}
\evid{23_dataflow_job.png}{Job \cod{lumina-fraude-stream-v3} (corrida 3) con workers sin IP pública en la subred \cod{lumina-datos-us-central1}. Lo lanzamos solo con \cod{--subnetwork}; en ese caso la API deja \cod{network: default} en el descriptor, pero la red la define la subred, que solo existe en \cod{lumina-vpc} (figura~\ref{fig:red})}{dataflowjob}
\evid{consola/dataflow_grafo.png}{Grafo del job \cod{lumina-fraude-stream-v3}: lectura de Pub/Sub, parseo con salida a DLQ y etapas con estado por cliente. \cod{ViajeImposible} concentra el 37 \% del tiempo de procesamiento}{c-grafo}

\subsection{Pruebas de carga}

\evid{26a_update_options.png}{Ajuste en caliente del rango de workers con \cod{update-options}, el mismo mecanismo que usa el DAG de precalentamiento}{updateoptions}
\evid{26b_carga_corrida1.png}{Corrida 1: eventos por minuto y latencia, con el aviso de lecturas de estado lentas}{corrida1}
\evid{27_pubsub_backlog_corrida1.png}{Corrida 1: backlog de la suscripción y antigüedad del mensaje más viejo}{backlog}
\evid{26c_carga_corrida2.png}{Corrida 2: eventos por minuto y uso de CPU del worker}{corrida2}
\evid{26d_carga_corrida3.png}{Corrida 3: eventos por minuto con la escritura a BigQuery en lotes. El primer minuto es la prueba de calentamiento a 5 ev/s, con p95 de 2,4 s}{corrida3}
\evid{consola/dataflow_metricas.png}{Consola, referencial: métricas del job \cod{v3}: un solo worker (máximo 1) y la CPU sube a 97 \% cuando entra la carga}{c-metricas}

\subsection{Gold y BigQuery ML}

\evid{29a_gold.png}{Creación de los hechos y vistas de la capa Gold}{goldcrea}
\evid{29_gold.png}{Objetos de Gold y KPI operativos por origen (batch y streaming)}{gold}
\evid{29b_reglas.png}{Desempeño de las reglas del pipeline e ingreso por producto}{reglas}
\evid{30a_bqml.png}{Entrenamiento de los modelos de fraude y liquidez}{bqml}
\evid{30_bqml_eval.png}{Evaluación del modelo de fraude y variables más importantes}{bqmleval}
\evid{30b_liquidez.png}{Error del pronóstico por día y plan de liquidez T+1 por distrito con la vista \cod{v\_plan\_liquidez\_t1} corregida}{liquidez}
\evid{30d_bqml_sin_fuga.png}{Modelo de fraude con y sin las variables posteriores a la transacción, y variables más importantes del modelo sin fuga}{sinfuga}
\evid{30c_mape.png}{MAPE del pronóstico por distrito, con y sin el último día de PaySim}{mape}

\subsection{API de datos en Cloud Run}

\evid{31_cloud_run_deploy.png}{Despliegue de la API con concurrencia 20 y de 0 a 3 instancias}{deploy}
\evid{31b_cloud_run_config.png}{Configuración desplegada: máximo de 3 instancias, concurrencia 20, 1 vCPU y 512 MiB. La URL \cod{lumina-api-gk5x4vgopq-uc.a.run.app} y la de la figura~\ref{fig:deploy}, basada en el número de proyecto, son dos alias que Cloud Run da a la misma revisión \cod{lumina-api-00001-6zb}}{cloudrun}
\evid{consola/cloud_run.png}{Métricas de \cod{lumina-api} en Cloud Run: solicitudes 2xx y 4xx, latencia e instancias que vuelven a cero cuando no hay tráfico}{c-cloudrun}
\evid{32_api_sin_token.png}{Sin credenciales la API responde 403}{apisintoken}
\evid{33_api.png}{Con token, la API devuelve la alerta más reciente y el plan de liquidez T+1 del distrito 1}{api}


% ====================================================================== 7
\section{Pendientes para la Entrega 3}

\begin{itemize}
\item Conectar Looker Studio a las vistas \cod{v\_kpi\_*} y armar el tablero de la demo.
\item Desplegar los dos DAG en Managed Service for Apache Airflow y mostrar una ejecución completa.
\item Llevar el modelo de fraude al pipeline (hoy decide con reglas) y desplegar el servicio de autorización.
\item Optimizar las etapas con estado del pipeline y repetir la prueba de carga hasta llegar a 500 eventos por segundo por worker.
\item Repetir la prueba con autoscaling activo para confirmar que el job escala de 4 a 55 workers sin perder latencia y que la capacidad crece de forma lineal.
\item Medir por separado el efecto de la VPC privada en la capacidad del worker (corrida de carga 2).
\item Repetir la corrida de carga con \cod{--max\_cache\_memory\_usage\_mb} más alto, como sugirió el runner, antes de concluir que el límite es solo la CPU.
\item Dar a la API una cuenta de servicio propia, con lectura solo sobre las vistas seudonimizadas, y dejar de devolver \cod{cliente\_id} en claro.
\item Separar \cod{sa-lumina-pipelines} en una cuenta por componente, con roles sobre el dataset o bucket que usa.
\item Extender la llave CMEK a Silver, Bronze y las zonas raw y curated.
\item Asignar los policy tags a las columnas sensibles (\cod{cliente\_id} y saldos como \cod{PII\_ALTA}; \cod{lat}, \cod{lon} y \cod{device\_id} como \cod{PII\_MEDIA}) y mostrarlo con \cod{bq show --schema}.
\item Lanzar Dataproc Serverless en \cod{lumina-vpc} con \cod{--subnet lumina-datos-us-central1}.
\item Crear el esquema de Pub/Sub de \cod{tx-eventos} y dejar en Dataflow solo las reglas de dominio.
\item Poner el balanceador con Cloud Armor delante de la API y cerrar el acceso directo a su URL.
\item Pasar la normalización de PaySim al ETL de Spark.
\item Preparar la demo de punta a punta: simulador, Pub/Sub, Dataflow, alerta, API y tablero.
\end{itemize}

\clearpage
\appendix
\section{Cambios en el repositorio}

\begingroup\small
\begin{xltabular}{\linewidth}{@{}P{6.2cm}L@{}}
\toprule
\textbf{Archivo} & \textbf{Cambio}\\
\midrule\endhead
\cod{doc/arquitectura/lumina.c4} & Modelo C4 en LikeC4; \cod{render.sh} exporta las vistas\\
\cod{doc/evidencias/} & Logs de cada comando y generador de capturas\\
\cod{src/infra/\allowbreak gobierno\_seguridad\_red.sh} & VPC, firewall, KMS, auditoría, policy tags y Knowledge Catalog\\
\cod{src/sql/05\_gobierno\_seguridad.sql} & Vistas por entidad y señal de fraude compartida\\
\cod{src/sql/06\_validacion\_capacidad.sql} & Análisis de la prueba de carga\\
\cod{src/orchestration/\allowbreak dag\_lumina\_precalentamiento.py} & Sube y baja el mínimo de workers en días de pago\\
\cod{src/sql/02\_dimensiones.sql}, \cod{03\_gold.sql} & Trends con Colombia como referencia y SEC con el año 2018\\
\cod{src/infra/setup.sh} & Rol \cod{pubsub.viewer} para la cuenta de los pipelines\\
\cod{src/dataflow/\allowbreak fraud\_streaming\_pipeline.py} & Escritura a BigQuery en lotes y parámetros de ejecución medidos\\
\cod{src/sql/04\_bqml.sql} & Modelo de fraude sin variables posteriores a la transacción\\
\cod{src/infra/registrar\_linaje.py} & Registra el linaje de los pasos de BigQuery y del streaming\\
\bottomrule
\end{xltabular}
\endgroup

\end{document}
